\subsection{Marco internacional aplicable al proyecto}\label{marco-internacional-aplicable-al-proyecto}

\subsection{Apertura}\label{apertura}

El proyecto de Transportes Curimón S.A. se ejecuta en Chile, pero su cadena de obligaciones no termina en la frontera. Tres situaciones lo internacionalizan. La primera es la ruta Antofagasta a Puerto Montt con cruce a Mendoza, que mueve datos de conductores y de carga a un tercer país. La segunda es la elección de Azure Chile Central como región primaria y Azure East US 2 como réplica, que traslada copias fuera del territorio nacional. La tercera es que el mandante exige controles de las familias ISO/IEC 27001 y 27002 en el requisito RT-11.05, es decir, importa un estándar internacional al pliego mismo de la licitación.

De ahí que el marco europeo no aparezca aquí como referencia académica sino por dos vías concretas. Una es directa, cuando la norma europea se aplica por sí misma. La otra es indirecta, cuando la norma europea opera como patrón de diseño de la regulación chilena, como ocurre con la Ley 21.719 respecto del RGPD.

El estado de vigencia de cada una de estas normas, con su fuente oficial y su fecha de consulta, se documenta en el \textbf{Anexo B}. De esa verificación se desprende el hallazgo que ordena el resto del capítulo: \textbf{ninguna de las normas europeas obliga hoy a Curimón por sí misma}, pero todas operan como patrón de diseño de la regulación chilena o como exigencia contractual probable de sus clientes.

\subsection{Análisis de la asimetría normativa}\label{anuxe1lisis-de-la-asimetruxeda-normativa}

La falta de cláusulas del modelo chileno crea un riesgo critico a la continuidad operativa de Curimon, lo que puede resultar en plazos extendidos ya que los contratos ya firmados con el proveedor de nube y con los transportistas externos deben rehacerse,aumento de costo al tener que redactar clausulas propias y someterlas a revisión legal y problemas de arquitectura al tener plataformas internacionales y nacionales como Azure Chile Central. Mientras no existan cláusulas modelo nacionales, la propuesta incorpora cláusulas contractuales basadas en el capítulo V del RGPD como diseño de preferencia, declarando expresamente que se trata de un instrumento provisorio, y reserva en el presupuesto las horas de revisión legal para su reemplazo.

\subsection{Normas técnicas certificables}\label{normas-tuxe9cnicas-certificables}

\subsection{Las cuatro de la ficha}\label{las-cuatro-de-la-ficha}

{\def\LTcaptype{none} % do not increment counter
\begin{longtable}[]{@{}
  >{\raggedright\arraybackslash}p{(\linewidth - 6\tabcolsep) * \real{0.2353}}
  >{\raggedright\arraybackslash}p{(\linewidth - 6\tabcolsep) * \real{0.2353}}
  >{\centering\arraybackslash}p{(\linewidth - 6\tabcolsep) * \real{0.2941}}
  >{\raggedright\arraybackslash}p{(\linewidth - 6\tabcolsep) * \real{0.2353}}@{}}
\toprule\noalign{}
\begin{minipage}[b]{\linewidth}\raggedright
Norma
\end{minipage} & \begin{minipage}[b]{\linewidth}\raggedright
Objeto
\end{minipage} & \begin{minipage}[b]{\linewidth}\centering
Certificable
\end{minipage} & \begin{minipage}[b]{\linewidth}\raggedright
Nota de vigencia
\end{minipage} \\
\midrule\noalign{}
\endhead
\bottomrule\noalign{}
\endlastfoot
ISO/IEC 27001:2022 & Sistema de gestión de seguridad de la información & Sí & El plazo de transición desde la edición 2013 venció en octubre de 2025. Toda certificación vigente está en la edición 2022 \\
ISO/IEC 27701:2025 & Sistema de gestión de privacidad (PIMS) & Sí & \textbf{Cambio estructural.} La edición 2025 la convierte en norma autónoma. Ya no requiere un SGSI 27001 previo como sí exigía la edición 2019 \\
ISO/IEC 42001:2023 & Sistema de gestión de inteligencia artificial & Sí & Vigente. Aplicable al modelo predictivo de fatiga y al bloqueo automático de despacho \\
NIST CSF 2.0 & Marco de gestión de ciberseguridad & No certificable & Publicado en 2024. Incorpora la función Govern. Se usa como marco de madurez, no como sello \\
\end{longtable}
}

\textbf{Por qué el cambio de la 27701 importa al costeo.} Bajo la edición 2019 la ruta de privacidad certificable obligaba a pagar antes una certificación 27001 completa. Bajo la edición 2025 la privacidad puede certificarse por separado. Eso abre una alternativa de secuenciación que el modelo de costos de P4 debe poder evaluar, en lugar de asumir una sola ruta. La consecuencia práctica es que la partida de certificación deja de ser un número único y pasa a ser una decisión de camino.

\subsection{Dos marcos adicionales de aporte propio}\label{dos-marcos-adicionales-de-aporte-propio}

La ficha cierra su lista con «y otros que el grupo debe identificar». El punto 4 de las Indicaciones obliga a agregar al menos dos alternativas y a justificar qué aportan. El grupo incorpora las siguientes.

\textbf{a) Ley argentina 25.326 de Protección de los Datos Personales.} Justificación de pertinencia: el Caso 10 contempla cruce fronterizo a Mendoza. Ningún otro marco de la ficha cubre el tratamiento de datos que ocurre del lado argentino de la operación. Es el único marco adicional que es obligatorio y no voluntario para este caso concreto.

\textbf{b) SOC 2 Type II (AICPA, TSC 2017 con revisiones posteriores).} Justificación de pertinencia: es el informe de aseguramiento que con más frecuencia exigen los clientes corporativos en contratos de servicios tecnológicos. Curimón atiende a 84 empresas cliente. A diferencia de ISO 27001, no certifica un sistema de gestión sino la operación efectiva de controles durante un periodo observado, lo que lo hace complementario y no sustituto.

Ninguna de las certificadoras que operan en Chile publica precio de lista para la auditoría de certificación: el régimen aplicable es \textbf{«solo por cotización»} conforme al punto 5 de las Indicaciones, y así se declara. Las cuatro componentes separables del costo ---compra del texto de la norma, consultoría de implantación, auditoría de certificación y mantención--- se incorporan al modelo económico del capítulo 3.
